\documentclass[11pt,a4paper]{article}
\usepackage[T1]{fontenc}
\usepackage[utf8]{inputenc}
\usepackage{lmodern,amsmath,amssymb}
\usepackage[margin=25mm]{geometry}
\usepackage{longtable,booktabs,array,calc}
\usepackage[hidelinks]{hyperref}
\hypersetup{pdftitle={An explicit strong-coupling gap for the Wilson transfer matrix: $SU(3)$ is gapped for $g^2},pdfauthor={navstokgap research working notes}}
\newcommand{\tightlist}{\setlength{\itemsep}{0pt}\setlength{\parskip}{0pt}}
\setlength{\emergencystretch}{3em}
\setcounter{secnumdepth}{0}
\begin{document}
\hypertarget{an-explicit-strong-coupling-gap-for-the-wilson-transfer-matrix-su3-is-gapped-for-g2ge176-with-deltagehbar-ca4logg2176}{%
\section{\texorpdfstring{An explicit strong-coupling gap for the Wilson
transfer matrix: \(SU(3)\) is gapped for \(g^2\ge176\), with
\(\Delta\ge(\hbar c/a)\,4\log(g^2/176)\)}{An explicit strong-coupling gap for the Wilson transfer matrix: SU(3) is gapped for g\^{}2\textbackslash ge176, with \textbackslash Delta\textbackslash ge(\textbackslash hbar c/a)\textbackslash,4\textbackslash log(g\^{}2/176)}}\label{an-explicit-strong-coupling-gap-for-the-wilson-transfer-matrix-su3-is-gapped-for-g2ge176-with-deltagehbar-ca4logg2176}}

The threshold that Yarotsky's theorem supplies for the Kogut--Susskind
Hamiltonian is \(g_0^2\sim10^{101}\)
(\href{strong-coupling-threshold-explicit.md}{threshold note}). The same
statement for the other standard lattice Hamiltonian, the logarithm of
the Wilson transfer matrix, has an explicit and usable threshold,
because the Euclidean character expansion is a polymer gas of closed
plaquette surfaces with no time discretization to pay for. With the
plaquette activity \(\rho\) of the fundamental representation, which is
\(1/g^2\) for \(SU(3)\) to leading order, the Kotecký--Preiss criterion
holds for
\[20e^2\rho\ \le\ 0.84,\qquad\text{that is}\qquad \rho\le5.7\times10^{-3},\]
and the truncated correlations of local observables then decay with rate
at least \(4\log(1/(176\rho))\) per lattice unit, because the smallest
closed surface joining two distant plaquettes is a tube with four
plaquettes per unit length. Through the transfer matrix
\(\mathcal T=e^{-aH_W/\hbar c}\), self-adjoint and positive by link
reflection positivity, the decay rate is a lower bound on the gap of
\(H_W\) on the cyclic subspace of local observables, so for \(SU(3)\)
\[\Delta_W\ \ge\ \frac{\hbar c}{a}\,4\log\frac{g^2}{176}\qquad(g^2\ge176),\]
uniformly in the volume. The fully rigorous version, which sums all
representations and both orientations with the bound
\(e^{\beta_W}-1\simeq6/g^2\), gives the same form with \(1056\) in place
of \(176\). The entropy constant \(20e\) for plaquette animals is the
only place with room, and a careful count would lower \(176\) to about
\(10^2\). The gap of \(H_W\) grows like \(\log g^2\) where that of
\(H_{\rm KS}\) grows like \(g^2\), the two Hamiltonians being different
regularizations that agree only in the continuum limit. Constants
explicit; nothing promoted.

\hypertarget{the-polymer-gas}{%
\subsection{1. The polymer gas}\label{the-polymer-gas}}

Write the Wilson weight of a plaquette in characters,
\[e^{\frac{\beta_W}{N}\operatorname{Re}\operatorname{tr}U_p}
=c_0(\beta_W)\Big[1+\sum_{r\ne0}d_r\frac{c_r}{c_0}\chi_r(U_p)\Big],
\qquad c_r=\frac1{d_r}\int e^{\frac{\beta_W}{N}\operatorname{Re}\operatorname{tr}U}\,\overline{\chi_r(U)}\,dU ,\]
with \(\beta_W=2N/g^2\). Expanding the product over plaquettes and
integrating over links, a set of plaquettes with nontrivial
representations contributes only if the representations at every link
fuse to the trivial one, so every link is covered by at least two
plaquettes of the set: the surviving sets are \textbf{closed surfaces},
the smallest being the six faces of a cube. On a closed surface carrying
the fundamental on every plaquette, the link integrals
\(\int\chi_f(U_1g)\chi_f(g^{-1}U_2)dg=\chi_f(U_1U_2)/d_f\) and the
vertex closures combine to \(d_f^{\chi(P)}(c_f/c_0)^{|P|}\) with
\(\chi(P)\) the Euler characteristic, so the activity per plaquette
along a tube, where the vertex and link factors balance, is
\[\rho\ =\ d_f\,\frac{c_f}{c_0}\ =\ \frac{\beta_W}{2N}\big(1+O(\beta_W)\big)=\frac1{g^2}\big(1+O(g^{-2})\big)\qquad(SU(N)),\]
the textbook strong-coupling parameter. The polymers of the gas are the
connected closed surfaces, two polymers being incompatible when they
share a link, and
\(Z=c_0^{|\mathcal P|}\sum_{\text{compatible}}\prod w(\gamma)\).

\textbf{All representations, rigorously.} Expand
\(e^{x\operatorname{Re}\operatorname{tr}U}=\sum_kx^k(\operatorname{Re}\operatorname{tr}U)^k/k!\)
with \(x=\beta_W/N\), and
\((\operatorname{Re}\operatorname{tr}U)^k=2^{-k}\sum_j\binom kj(\operatorname{tr}U)^j(\operatorname{tr}U^\dagger)^{k-j}\).
The integral
\(\int(\operatorname{tr}U)^j(\operatorname{tr}U^\dagger)^{k-j}\overline{\chi_r}\,dU\)
is the multiplicity \(m_r(j,k-j)\) of \(r\) in
\(f^{\otimes j}\otimes\bar f^{\otimes(k-j)}\), so
\(d_rc_r=\sum_k\frac{x^k}{k!}2^{-k}\sum_j\binom kjm_r(j,k-j)\ge0\), and
\(\sum_rm_r(j,k-j)\le\dim\big(f^{\otimes j}\otimes\bar f^{\otimes(k-j)}\big)=N^k\).
Hence
\[\sum_{r\ne0}d_rc_r\ \le\ \sum_{k\ge1}\frac{x^k}{k!}2^{-k}\cdot2^kN^k=e^{Nx}-1=e^{\beta_W}-1 ,\]
and Jensen's inequality with
\(\int\operatorname{Re}\operatorname{tr}U\,dU=0\) gives
\(c_0=\int e^{x\operatorname{Re}\operatorname{tr}U}dU\ge1\). Therefore
\[\sum_{r\ne0}d_r\frac{c_r}{c_0}\ \le\ e^{\beta_W}-1=e^{2N/g^2}-1\ \simeq\ \frac{2N}{g^2}=\frac{6}{g^2}\quad(SU(3)),\]
which is the leading activity \(1/g^2\) of the fundamental, its
conjugate, and the higher representations together, with no spurious
factor. The sum over representations also covers the branch lines where
three fundamentals meet at a link, allowed for \(SU(3)\) since
\(f^{\otimes3}\ni1\). Everything below is stated for a per-plaquette
activity \(\rho\); the two readings differ only in what \(\rho\) is.

\hypertarget{convergence-the-koteckuxfdpreiss-criterion-with-numbers}{%
\subsection{2. Convergence: the Kotecký--Preiss criterion with
numbers}\label{convergence-the-koteckuxfdpreiss-criterion-with-numbers}}

\emph{Adjacency.} In four dimensions a link lies in six plaquettes, so a
plaquette shares a link with \(4\times5=20\) others, and the number of
connected sets of \(n\) plaquettes containing a given one is at most
\((20e)^{n-1}\).

\emph{Incompatible polymers.} A polymer \(\gamma'\) incompatible with
\(\gamma\) shares a link with it; \(\gamma\) has at most \(4|\gamma|\)
links, each in six plaquettes, so the number of \(\gamma'\) of size
\(n\) is at most \(24|\gamma|(20e)^{n-1}\), and \(n\ge6\).

\emph{Criterion.} With \(a(\gamma)=|\gamma|\) and
\(|w(\gamma')|\le\rho^{|\gamma'|}\), the Kotecký--Preiss condition
\(\sum_{\gamma'\not\sim\gamma}|w(\gamma')|e^{a(\gamma')}\le a(\gamma)\)
follows from \[\frac{24}{20e}\sum_{n\ge6}\big(20e^2\rho\big)^n\ \le\ 1
\qquad\Longleftarrow\qquad x\equiv20e^2\rho\le0.84 ,\] since
\(0.4415\,x^6/(1-x)\le1\) there. So the expansion converges, uniformly
in the volume, for \[\rho\ \le\ \frac{0.84}{20e^2}=5.7\times10^{-3}.\]

\hypertarget{decay-of-correlations-and-the-gap}{%
\subsection{3. Decay of correlations and the
gap}\label{decay-of-correlations-and-the-gap}}

Running the criterion with \(a(\gamma)=(1+\delta)|\gamma|\) instead, it
holds as long as \(xe^{\delta}\le0.84\), that is for
\(\delta\le\log(0.84/x)=\log(1/(176\rho))\), and then the cluster sums
carry a factor \(e^{-\delta|X|}\). A cluster connecting two plaquettes
at distance \(T\) contains a closed connected surface through both,
whose size is at least that of a tube, \(4T\) plaquettes, so the
truncated two-point function of any two local observables at separation
\(T\) is bounded by \(C\,e^{-4\delta T}\):
\[m\,a\ \ge\ 4\log\frac{1}{176\rho}.\]

\textbf{Transfer.} The Wilson action is link-reflection positive, so the
transfer matrix \(\mathcal T\) is self-adjoint and positive (Lüscher,
Commun. Math. Phys. 54 (1977) 283; Osterwalder and Seiler, Ann. Phys.
110 (1978) 440; metadata level), and \(H_W=-(\hbar c/a)\log\mathcal T\).
For a local observable \(A\) orthogonal to the vacuum,
\(\langle A,\mathcal T^nA\rangle\le C_Ae^{-man}\) for all \(n\) forces
the spectral measure of \(A\) to vanish above \(e^{-ma}\), so on the
cyclic subspace generated by local observables, which is the whole space
of the reconstructed theory,
\[\Delta_W\ \ge\ \hbar c\,m\ \ge\ \frac{\hbar c}{a}\,4\log\frac{1}{176\rho}.\]

\textbf{Numbers for \(SU(3)\).}

\begin{longtable}[]{@{}
  >{\raggedright\arraybackslash}p{(\columnwidth - 4\tabcolsep) * \real{0.3333}}
  >{\raggedright\arraybackslash}p{(\columnwidth - 4\tabcolsep) * \real{0.3333}}
  >{\raggedright\arraybackslash}p{(\columnwidth - 4\tabcolsep) * \real{0.3333}}@{}}
\toprule\noalign{}
\begin{minipage}[b]{\linewidth}\raggedright
reading of \(\rho\)
\end{minipage} & \begin{minipage}[b]{\linewidth}\raggedright
threshold \(\rho\le5.7\times10^{-3}\)
\end{minipage} & \begin{minipage}[b]{\linewidth}\raggedright
gap
\end{minipage} \\
\midrule\noalign{}
\endhead
\bottomrule\noalign{}
\endlastfoot
leading, \(\rho=1/g^2\) & \(g^2\ge176\) &
\(\Delta_W\ge(\hbar c/a)\,4\log(g^2/176)\) \\
all representations, \(\rho=e^{6/g^2}-1\) & \(g^2\ge1056\) &
\(\Delta_W\ge(\hbar c/a)\,4\log(g^2/1056)\) \\
\end{longtable}

The first line uses the fundamental alone; the second sums every
representation and both orientations and is rigorous as written. The
factor \(6\) between them is the two orientations, each with activity
\(1/g^2\), and the higher representations, so the rigorous line is close
to the truth for a surface that may carry either orientation on each
plaquette; a surface with a fixed global orientation would be counted by
the first line.

\hypertarget{relation-to-the-kogutsusskind-statement}{%
\subsection{4. Relation to the Kogut--Susskind
statement}\label{relation-to-the-kogutsusskind-statement}}

\(H_W\) and \(H_{\rm KS}\) are different regularizations of the same
continuum Hamiltonian. At strong coupling their gaps differ in kind:
\(H_{\rm KS}\) has an electric term of order \(g^2\hbar c/a\) and the
flux loop costs \((8/3)g^2\hbar c/a\) for \(SU(3)\); \(H_W\) has both
electric and magnetic content in \(-\log\mathcal T\), and the flux loop
costs \(4\log(g^2/\cdot)\,\hbar c/a\), the logarithm of the tube
activity. The Jaffe--Witten statement asks for the continuum theory and
is indifferent to the regularization, so either Hamiltonian serves for
T2, and the one with the explicit threshold is \(H_W\). The anisotropic
limit \(a_t\to0\) that turns \(\mathcal T\) into
\(e^{-a_tH_{\rm KS}/\hbar}\) sends the temporal plaquette activity to
one, outside the expansion, which is why the Kogut--Susskind threshold
needs a genuinely Hamiltonian expansion and why Yarotsky's time
discretization pays what it pays.

\hypertarget{consequence-for-state}{%
\subsection{5. Consequence for STATE}\label{consequence-for-state}}

The strong-coupling boundary of the intermediate region is now a number:
for the Wilson transfer matrix and \(SU(3)\), the theory is gapped
uniformly in the volume for \(g^2\ge176\) (leading activity) or
\(g^2\ge1056\) (all representations, rigorous), with the gap explicit.
The only constant with room is the plaquette-animal entropy \(20e\). On
the weak side the boundary is \(1/g^2\gtrsim10^2\) with the crude window
of \href{large-field-operator-inequality.md}{the operator-inequality
note}, and the intermediate region between \(g^2\sim10^{-2}\) and
\(g^2\sim2\times10^2\) is where the gap forms.
\end{document}
